% =========================================================
% Reflective activity -- Time Series Analysis module
% Author : Aymen Ben Brik (aymen.benbrik@esprit.tn)
% Esprit School of Business
% =========================================================
\documentclass[a4paper,11pt]{article}

\newcommand{\TPnumero}{Reflective activity}
\newcommand{\macrosPath}{../preamble/macros.tex}
% =========================================================
% Tutorial / lab preamble — Time Series Analysis module
% article class + fancyhdr + tcolorbox + listings (English)
% Author : Aymen Ben Brik (aymen.benbrik@esprit.tn)
% =========================================================

\usepackage[utf8]{inputenc}
\usepackage[T1]{fontenc}
\usepackage[english]{babel}
\usepackage{lmodern}
\usepackage{amsmath, amssymb, amsthm, mathtools, bm}
\usepackage{graphicx}
\usepackage{booktabs}
\usepackage{array}
\usepackage{multirow}
\usepackage{colortbl}
\usepackage{xcolor}
\usepackage{tikz}
\usetikzlibrary{positioning, arrows.meta, calc, shapes, fit, backgrounds, matrix}
\usepackage{listings}
\usepackage[most]{tcolorbox}
\usepackage{geometry}
\geometry{a4paper, margin=2cm, headheight=15pt}
\usepackage{fancyhdr}
\usepackage[hidelinks]{hyperref}
\usepackage{enumitem}

% --- Colours ---
\definecolor{espritBlue}{RGB}{30, 60, 110}
\definecolor{espritAccent}{RGB}{200, 80, 30}
\definecolor{boxEx}{RGB}{30, 60, 110}
\definecolor{boxQ}{RGB}{180, 120, 0}
\definecolor{boxC}{RGB}{30, 100, 60}
\definecolor{boxAtt}{RGB}{200, 80, 30}

% --- Listings ---
\definecolor{codebg}{RGB}{248,248,248}
\definecolor{codeframe}{RGB}{220,220,220}
\definecolor{keyword}{RGB}{0,82,155}
\definecolor{comment}{RGB}{88,110,117}
\definecolor{string}{RGB}{133,153,0}

\lstdefinestyle{pythonstyle}{
  language=Python,
  basicstyle=\ttfamily\small,
  keywordstyle=\color{keyword}\bfseries,
  commentstyle=\color{comment}\itshape,
  stringstyle=\color{string},
  backgroundcolor=\color{codebg},
  frame=single,
  rulecolor=\color{codeframe},
  frameround=tttt,
  breaklines=true,
  showstringspaces=false,
  tabsize=2
}
\lstdefinestyle{Rstyle}{
  language=R,
  basicstyle=\ttfamily\small,
  keywordstyle=\color{keyword}\bfseries,
  commentstyle=\color{comment}\itshape,
  stringstyle=\color{string},
  backgroundcolor=\color{codebg},
  frame=single,
  rulecolor=\color{codeframe},
  frameround=tttt,
  breaklines=true,
  showstringspaces=false,
  tabsize=2
}
\lstset{style=pythonstyle}

% --- Common macros (configurable path) ---
\providecommand{\macrosPath}{../../preamble/macros.tex}
\input{\macrosPath}

% --- Exercise counter ---
\newcounter{excounter}
\renewcommand{\theexcounter}{\arabic{excounter}}

\newtcolorbox[use counter=excounter]{exercise}[1][]{
  enhanced, breakable,
  colback=boxEx!5, colframe=boxEx, fonttitle=\bfseries,
  title={Exercise~\theexcounter\ifstrempty{#1}{}{ -- #1}},
  arc=2pt, boxrule=0.6pt
}
\newtcolorbox{question}[1][]{
  enhanced, breakable,
  colback=boxQ!5, colframe=boxQ, fonttitle=\bfseries,
  title={Question\ifstrempty{#1}{}{ -- #1}},
  arc=2pt, boxrule=0.5pt
}
\newtcolorbox{solution}[1][]{
  enhanced, breakable,
  colback=boxC!5, colframe=boxC, fonttitle=\bfseries,
  title={Solution\ifstrempty{#1}{}{ -- #1}},
  arc=2pt, boxrule=0.5pt
}
\newtcolorbox{warning}[1][]{
  enhanced, breakable,
  colback=boxAtt!5, colframe=boxAtt, fonttitle=\bfseries,
  title={Warning\ifstrempty{#1}{}{ -- #1}},
  arc=2pt, boxrule=0.5pt
}

% --- Headers / footers ---
\pagestyle{fancy}
\fancyhf{}
\fancyhead[L]{\textsc{\moduleNom} -- \auteurNom}
\fancyhead[R]{\TPnumero}
\fancyfoot[C]{\thepage}
\fancyfoot[L]{\auteurInstitution}
\fancyfoot[R]{\auteurEmail}
\renewcommand{\headrulewidth}{0.4pt}
\renewcommand{\footrulewidth}{0.2pt}

% --- Placeholder ---
\providecommand{\TPnumero}{Lab}


\title{End-of-module reflective activity\\
\large \moduleNom\ -- guided self-assessment}
\author{\auteurNom\\\auteurEmail\\\auteurInstitution}
\date{\anneeUniversitaire}

\begin{document}
\maketitle

\section*{Instructions}

This \emph{ungraded} activity is intended to help you:
\begin{itemize}
  \item identify which notions of the module you have mastered and
        which ones still need work;
  \item formalise what you retain from the six chapters;
  \item plan your revision for the final exam and the integrative
        project.
\end{itemize}

\noindent
Take about \textbf{45 minutes} to answer in writing (handwritten or
in LaTeX). Be honest: this document will not be collected and the
benefit lies in the metacognitive effort you invest.

% =========================================================
\begin{question}[1.~Self-positioning by chapter]
For each chapter, give yourself a \emph{perceived mastery} score
from $1$ (very fuzzy) to $5$ (I could explain it to a classmate)
and justify in one sentence.

\medskip
\begin{tabular}{lcl}
\toprule
\textbf{Chapter} & \textbf{Score 1--5} & \textbf{Justification (one sentence)} \\
\midrule
Ch.1 -- Basic statistics       & \rule{1cm}{0.4pt} & \dotfill \\
Ch.2 -- Decomposition          & \rule{1cm}{0.4pt} & \dotfill \\
Ch.3 -- Stationarity (ACF/PACF) & \rule{1cm}{0.4pt} & \dotfill \\
Ch.4 -- Non-stationarity (ADF/KPSS) & \rule{1cm}{0.4pt} & \dotfill \\
Ch.5 -- ARMA models            & \rule{1cm}{0.4pt} & \dotfill \\
Ch.6 -- ARCH/GARCH             & \rule{1cm}{0.4pt} & \dotfill \\
\bottomrule
\end{tabular}
\end{question}

% =========================================================
\begin{question}[2.~The hardest notion and why]
Identify the notion of the module that you found most
\emph{difficult}. Describe:
\begin{itemize}
  \item what makes it difficult for you (algebra, intuition,
        probability, the link with the code\dots);
  \item one concrete strategy to work on it next week (revisit a
        specific example, redo a lab exercise, draw a diagram,
        re-derive a formula by hand\dots).
\end{itemize}
\end{question}

% =========================================================
\begin{question}[3.~Sketch the Box--Jenkins workflow from memory]
Without looking at the lecture notes, sketch the Box--Jenkins
workflow from a raw series $X_t$ all the way to a validated ARIMA
forecast. Make sure your diagram includes:
\[
  \text{raw } X_t
  \;\to\; \text{stationarise}
  \;\to\; \text{identify }(p, q)
  \;\to\; \text{estimate}
  \;\to\; \text{validate}
  \;\to\; \text{forecast.}
\]
For each arrow, write the \emph{tool} you would use (e.g.\ ADF,
ACF, MLE, Ljung--Box, MA$(\infty)$ variance). Then compare with the
\emph{Box--Jenkins identification} section of Chapter~5 in the
lecture notes and correct any missing or misplaced step.
\end{question}

% =========================================================
\begin{question}[4.~Quick mental computations]
\textbf{(a)} For $X_t = 0.7\, X_{t-1} + \varepsilon_t$ with
$\varepsilon_t \iid \mathcal N(0, 1)$, write the values of
$\rho(1), \rho(2), \rho(3)$. Is the PACF of this AR(1) larger than
the ACF at lag~$2$? Why?

\medskip
\textbf{(b)} For an MA(1) process $X_t = \varepsilon_t + \theta
\varepsilon_{t-1}$ with $\theta = 0.5$, what is $\rho(1)$? At
which value of $\theta$ does the maximum of $|\rho(1)|$ occur?

\medskip
\textbf{(c)} For a GARCH(1,1) with $(\alpha_0, \alpha_1, \beta_1) =
(0.05, 0.10, 0.85)$, compute the unconditional variance
$\sigma^2$ and the half-life of a volatility shock (number of days
$\tau$ such that $(\alpha_1 + \beta_1)^{\tau} = 1/2$).

\medskip
\textbf{(d)} You run an ADF test and obtain stat $= -1.5$, while
KPSS gives stat $= 1.8$. Both at $5\%$ levels. In which cell of
the ADF/KPSS combination matrix does this land? What action would
you take?
\end{question}

% =========================================================
\begin{question}[5.~From maths to code]
Pick \emph{one} mathematical formula from the module and explain:
\begin{itemize}
  \item how you \emph{implemented} it in a lab (in NumPy / pandas /
        statsmodels / arch);
  \item one pitfall or bug you encountered (shape mismatch,
        off-by-one indexing, time-zone, log-difference vs simple
        difference, $\sigma^2$ versus $\sigma$, units (\% vs.\
        fraction)\dots);
  \item how you fixed it.
\end{itemize}
This question exercises the \emph{maths $\leftrightarrow$ code}
translation skill, which is one of the key competences of the
module.
\end{question}

% =========================================================
\begin{question}[6.~The thread that ties the six chapters]
In $5$--$10$ lines, explain why these six chapters are gathered
into a single module. What is the \emph{common thread}? What would
be lost by studying any one of them in isolation?

\medskip
\textit{Hint.} Think of the role of the residual $\widehat
\varepsilon_t$, which:
\begin{enumerate}
  \item appears in Chapter~2 as the residual of decomposition;
  \item is the object whose stationarity is tested in Chapter~3
        (ACF, Ljung--Box);
  \item is the object the unit-root tests are run on in Chapter~4;
  \item is what an ARMA model in Chapter~5 tries to make white;
  \item is the object whose conditional variance is modelled in
        Chapter~6.
\end{enumerate}
\end{question}

% =========================================================
\begin{question}[7.~Connecting to the integrative project]
Look back at the seven parts (A--G) of the final integrative
project (\texttt{projet-final/sujet.pdf}). For each part, write the
chapter it builds on:

\medskip
\begin{tabular}{ll}
\toprule
\textbf{Project part} & \textbf{Chapter mobilised} \\
\midrule
A -- Descriptive statistics + JB test & \dotfill \\
B -- Decomposition + linear trend     & \dotfill \\
C -- ACF/PACF + Ljung--Box            & \dotfill \\
D -- ADF/KPSS                         & \dotfill \\
E -- ARMA selection                   & \dotfill \\
F -- GARCH$(1, 1)$ + VaR              & \dotfill \\
G -- Kupiec back-test (synthesis)     & \dotfill \\
\bottomrule
\end{tabular}

\medskip
What does the back-test in Part~G say about \emph{what we have
not modelled}? In other words: even a passing Kupiec test is not a
proof of model correctness -- can you identify two limits of the
joint ARMA + GARCH model that the back-test does not see?
\end{question}

% =========================================================
\begin{question}[8.~Personal action plan]
List \emph{three} concrete actions you will take \emph{this week}
to consolidate your learning of the module:
\begin{enumerate}
  \item \dotfill
  \item \dotfill
  \item \dotfill
\end{enumerate}
A concrete action has an action verb (\emph{revisit}, \emph{code},
\emph{redo}, \emph{sketch}\dots), a precise object, and a
\emph{timing} (when will you do it?).
\end{question}

\bigskip
\noindent\textit{Good thinking. This activity will not be
collected: it is for you.}\par\medskip
\hfill\auteurNom\par
\hfill\auteurEmail\par

\end{document}
